\chapter{Memory Model and Synchronization}
\label{chap:memorymodel}

Chapter~\ref{chap:concurrency} lists the C++ APIs (\texttt{std::atomic}, mutexes, 
\texttt{memory\char`_order} enumerators). This chapter explains the \highlight{rules}: when two 
threads may see each other's writes, which orders are enough for a handoff, and how to 
build a small \highlight{lock-free SPSC} queue used on HFT hot paths, then two 
classic CAS structures (Treiber stack, Michael--Scott queue) that interviewers 
often expect next.

\section{C++ Memory Model}
\label{sec:cpp-memory-model}
The C++ memory model defines which program executions are allowed when threads share memory. 
Compilers and CPUs may \highlight{reorder} loads and stores unless synchronisation forbids it. 
Portable concurrent code relies on atomics, mutexes, or other synchronising operations --- 
not on ``it works on my machine.''

\subsection{Data Race}
\label{subsec:data-race}
A \highlight{data race} occurs when two threads access the \highlight{same memory location}, 
at least one access is a \highlight{write}, and the accesses are \highlight{not ordered} by 
the memory model (no mutex, no atomic with adequate ordering, no other happens-before edge).
\newline
In C++, a data race is \highlightbf{undefined behavior}. Fix it with atomics, locks, or by 
giving each thread private data.

\begin{lstlisting}[language=C++]
#include <thread>

int g = 0;  // shared, non-atomic

void bad()
{
    std::thread a([] { g = 1; });
    std::thread b([] { g = 2; });  // data race with a
    a.join();
    b.join();
}
\end{lstlisting}

\subsection{Race Condition}
\label{subsec:race-condition}
A \highlight{race condition} is a broader engineering term: the program's correctness depends on 
\highlight{timing} or interleaving. It may or may not be a formal data race.
\newline
Example: both threads use a mutex (no data race) but the business logic still breaks if 
updates run in the wrong order.
\par\noindent\textbf{Interview tip:} say \highlight{data race} when you mean UB; 
say \highlight{race condition} for logic/timing bugs.

\subsection{Happens-Before}
\label{subsec:happens-before}
\highlight{Happens-before} is the relation that makes writes in one thread \highlight{visible} 
to another. If A happens-before B, then B is guaranteed to see the side effects of A 
(for the relevant memory).
\newline
Typical ways to create happens-before between threads:
\begin{itemize}
    \item Unlock of a mutex happens-before a later lock of the \highlight{same} mutex.
    \item Atomic \texttt{release} store happens-before an atomic \texttt{acquire} load that 
          \highlight{reads the value written} by that store (on the same atomic object).
    \item Thread construction / \texttt{join} edges with the thread's entry and exit.
\end{itemize}
Without a happens-before edge, another thread may see stale or torn values.

\subsection{Sequential Consistency}
\label{subsec:sequential-consistency}
See Chapter~\ref{chap:concurrency}, Section~\ref{subsec:memory-order} for \texttt{memory\char`_order\char`_seq\char`_cst}. 
This subsection covers the formal model only.
\newline
Under \highlight{sequential consistency}, there is a single total order of all seq\_cst 
atomic operations that all threads agree on, and it is consistent with each thread's 
program order. It is the \highlight{easiest mental model} and the default for many atomic 
operations --- but often \highlight{stronger (and costlier)} than a producer/consumer handoff needs.

\section{Memory Orders}
\label{sec:memory-orders-theory}

\subsection{C++ Reference}
\label{subsec:memory-order-reference}
This subsection is a pointer, not a second API catalog. Chapter~\ref{chap:concurrency} 
owns the \texttt{memory\char`_order} / fence \highlight{surface} 
(Section~\ref{subsec:memory-sync-ordering}) and the when-to-pick bullets 
(Section~\ref{subsec:memory-order}).
\newline
Here the goal is \highlight{theory}: happens-before, release--acquire pairing, and 
why weakening is a trade between mental cost and hardware barriers. Read the 
Concurrency tables for names; stay in this chapter for proof sketches and lock-free 
structures.
\par\noindent\textbf{Interview tip:} say ``API in Concurrency, happens-before here'' --- 
do not recite every enumerator twice.

\subsection{Release-Acquire Pairing}
\label{subsec:release-acquire-pairing}
The standard pattern for publishing data from a producer to a consumer:
\begin{enumerate}
    \item Producer writes the payload (non-atomic or relaxed atomics).
    \item Producer performs a \highlight{release} store on a flag / index (``data is ready'').
    \item Consumer performs an \highlight{acquire} load on that flag / index.
    \item If the consumer sees the producer's value, it may safely read the payload.
\end{enumerate}
The release-acquire pair creates happens-before: all writes \highlight{before} the release 
become visible to the consumer \highlight{after} the matching acquire.

\begin{lstlisting}[language=C++]
#include <atomic>
#include <thread>

int payload = 0;
std::atomic<bool> ready{false};

void producer()
{
    payload = 42;                                  // (1)
    ready.store(true, std::memory_order_release);  // (2)
}

void consumer()
{
    while (!ready.load(std::memory_order_acquire)) {  // (3)
        // spin or wait
    }
    // (4) safe to read payload: sees 42
    int x = payload;
}
\end{lstlisting}

\noindent
\texttt{memory\char`_order\char`_acq\char`_rel} is for read-modify-write that both acquire and release 
(e.g.\ some lock implementations). \texttt{memory\char`_order\char`_relaxed} only guarantees 
atomicity of that location --- \highlight{no} synchronisation with other memory.

\subsection{When to Weaken Ordering}
\label{subsec:when-to-weaken-ordering}
Default atomics are often \texttt{seq\char`_cst}: one total order that is easy to reason about 
and \highlight{expensive} on weakly ordered CPUs (extra fences / store-buffer drains). 
Weakening is how you remove that barrier noise --- only after the publish/consume story is clear.
\begin{itemize}
    \item \highlightbf{seq\_cst} --- default / when unsure; simplest reasoning; may cost more on weak HW.
    \item \highlightbf{release / acquire} --- message passing, SPSC indices, ``publish flag'' patterns.
          \textbf{Noise removed:} full-system seq\_cst fences you did not need for a two-thread handoff.
    \item \highlightbf{relaxed} --- pure counters, statistics, when no other data depends on the atomic.
          \textbf{Noise removed:} acquire/release barriers on increments that never publish a payload.
    \item Prefer \highlight{correct first}; measure before weakening further on x86 (often strong in practice, 
          but portable code must still use the right orders --- ARM will expose the bug).
\end{itemize}
\par\noindent\textbf{Interview tip:} weaken only with a named pattern (counter vs publish/consume); 
never ``relaxed everywhere because x86.''

\section{Lock-Free Ring Buffer (SPSC)}
\label{sec:spsc-ring-buffer}
A \highlight{single-producer single-consumer} ring buffer moves messages between two threads 
without a mutex. It is a standard HFT building block (feed thread $\rightarrow$ strategy thread).
\newline
\textbf{Noise removed} versus a mutex queue: no kernel futex sleep/wake on the common path, 
no lock-word ping-pong, and a fixed preallocated buffer (no \texttt{new} per message). 
\textbf{Cost:} exactly one writer and one reader forever; overflow/empty must be handled 
explicitly (drop, block on a colder path, or size for burst).

\subsection{SPSC Model}
\label{subsec:spsc-model}
Exactly \highlight{one} thread pushes, exactly \highlight{one} thread pops. That restriction allows 
a simple design with two indices and no CAS loops for the common path: each index has a 
single writer, so updates need atomics for visibility, not contention resolution.
\newline
Multi-producer or multi-consumer queues need different algorithms 
(Sections~\ref{sec:treiber-stack}--\ref{sec:michael-scott-queue}) and are harder to get right. 
If your design needs MPMC on the hot path, revisit ownership first --- many HFT stacks 
shard into several SPSC rings instead.

\subsection{Head and Tail Indices}
\label{subsec:spsc-head-tail}
Use a fixed-capacity array. The producer advances a \highlight{tail} (write position); the consumer 
advances a \highlight{head} (read position). Indices are \texttt{std::atomic} (or indexed modulo capacity).
\newline
Only the producer writes \texttt{tail}; only the consumer writes \texttt{head}. Each thread may 
\highlight{read} the other's index to decide full/empty. That single-writer discipline is why 
SPSC avoids CAS on the steady path.

\subsection{Producer and Consumer Ordering}
\label{subsec:spsc-producer-consumer-ordering}
Ordering is the same publish/consume story as Section~\ref{subsec:release-acquire-pairing}:
\begin{itemize}
    \item Producer: write slot, then \texttt{tail.store(..., release)}.
    \item Consumer: \texttt{head}-side acquire load of \texttt{tail} (or acquire load when observing 
          new data), then read slot.
\end{itemize}
Release-acquire on the indices publishes the slot contents to the consumer. 
\textbf{Noise removed by getting this right:} seeing a new index while the payload store is 
still invisible (or reading a slot after a relaxed bump that never synchronized).

\subsection{Full and Empty Conditions}
\label{subsec:spsc-full-empty}
With capacity $N$ (often power of two):
\begin{itemize}
    \item \highlight{Empty} when \texttt{head == tail} (consumer has caught up).
    \item \highlight{Full} when \texttt{tail - head == N} (no free slot) --- exact formula depends on 
          whether you leave one slot empty or use an explicit size counter.
\end{itemize}
Never let producer and consumer update the \highlight{same} index. Leaving one slot empty 
is a common trick so full and empty are distinguishable without an extra atomic size.

\subsection{Cache-Line Layout}
\label{subsec:spsc-cache-line-layout}
See Section~\ref{subsec:false-sharing-solutions} for padding and alignment as a general technique. 
This subsection covers SPSC-specific head/tail placement only.
\newline
Put \texttt{head} and \texttt{tail} on \highlight{separate cache lines} (\texttt{alignas(64)}; 
see Section~\ref{subsec:false-sharing-solutions} for the portable constant) so the 
producer's stores to \texttt{tail} do not invalidate the consumer's line for \texttt{head} 
(false sharing).

\subsection{Implementation}
\label{subsec:spsc-implementation}
Sketch (power-of-two capacity, one empty slot to distinguish full/empty):

\begin{lstlisting}[language=C++]
#include <atomic>
#include <cstdint>
#include <vector>

template <typename T, std::size_t N>
class SpscRing {
    static_assert((N & (N - 1)) == 0, "N power of two");
    static_assert(N >= 2);

    alignas(64) std::atomic<std::size_t> head_{0};  // consumer
    alignas(64) std::atomic<std::size_t> tail_{0};  // producer
    T buf_[N];

public:
    bool try_push(const T& v)
    {
        const std::size_t t = tail_.load(std::memory_order_relaxed);
        const std::size_t next = (t + 1) & (N - 1);
        if (next == head_.load(std::memory_order_acquire)) {
            return false;  // full
        }
        buf_[t] = v;
        tail_.store(next, std::memory_order_release);
        return true;
    }

    bool try_pop(T& out)
    {
        const std::size_t h = head_.load(std::memory_order_relaxed);
        if (h == tail_.load(std::memory_order_acquire)) {
            return false;  // empty
        }
        out = buf_[h];
        head_.store((h + 1) & (N - 1), std::memory_order_release);
        return true;
    }
};
\end{lstlisting}

\par\noindent\textbf{Interview tip:} SPSC only; release/acquire on indices; cache-line separation; 
preallocate \texttt{buf\_} (no \texttt{new} on the hot path).

\section{Lock-Free Stack (Treiber)}
\label{sec:treiber-stack}
See~\cite{wiki-treiber-stack}. The \highlight{Treiber stack} is the usual first 
multi-producer / multi-consumer lock-free structure after SPSC: a linked list with a 
single atomic \texttt{head}, updated with CAS.
\newline
Any thread may push or pop. Progress of one operation does not need a mutex, but a 
slow thread can still starve if CAS keeps failing (lock-free $\neq$ wait-free).

\subsection{Push and Pop}
\label{subsec:treiber-push-pop}
Both operations race on a single atomic \texttt{head}. Progress is lock-free because 
some CAS eventually wins; an individual thread may still spin under contention 
(lock-free $\neq$ wait-free, and $\neq$ ``fast under load'').
\begin{itemize}
    \item \highlightbf{push:} allocate a node, set \texttt{node->next = head}, 
          \texttt{compare\char`_exchange} \texttt{head} from old head to \texttt{node}; 
          retry on failure.
    \item \highlightbf{pop:} load \texttt{head}; if null, empty; else CAS 
          \texttt{head} from \texttt{old} to \texttt{old->next}; retry on failure; 
          then read the value out of \texttt{old}.
\end{itemize}
Ordering sketch: \texttt{release} on successful push publish; \texttt{acquire} on 
pop's load/CAS so the node fields are visible.
\newline
\textbf{Noise removed} versus a mutex stack: no futex; failed CAS retries in user space. 
\textbf{Noise introduced:} allocator traffic per node, ABA, and reclamation --- why HFT 
hot paths usually prefer preallocated rings over unbounded Treiber stacks.

\subsection{Implementation Sketch}
\label{subsec:treiber-implementation}

\begin{lstlisting}[language=C++]
#include <atomic>
#include <utility>

template <typename T>
class TreiberStack {
    struct Node {
        T value;
        Node* next;
    };

    std::atomic<Node*> head_{nullptr};

public:
    void push(T v)
    {
        Node* n = new Node{std::move(v), nullptr};
        Node* old = head_.load(std::memory_order_relaxed);
        do {
            n->next = old;
        } while (!head_.compare_exchange_weak(
            old, n,
            std::memory_order_release,
            std::memory_order_relaxed));
    }

    bool try_pop(T& out)
    {
        Node* old = head_.load(std::memory_order_acquire);
        do {
            if (old == nullptr) {
                return false;
            }
        } while (!head_.compare_exchange_weak(
            old, old->next,
            std::memory_order_acquire,
            std::memory_order_acquire));

        out = std::move(old->value);
        // delete old;  // unsafe without a reclamation scheme (see below)
        return true;
    }
};
\end{lstlisting}

\subsection{ABA and Memory Reclamation}
\label{subsec:treiber-aba}
Two hard parts interviewers probe:
\begin{itemize}
    \item \highlightbf{ABA:} thread~A reads \texttt{head == X}; another thread pops 
          \texttt{X}, pushes other nodes, then pushes a recycled \texttt{X}; A's CAS 
          succeeds even though the list changed. Mitigations: tagged pointers 
          (version + pointer in a wide CAS), or never reuse addresses while a 
          lagged reader might still CAS them.
    \item \highlightbf{Safe reclamation:} you cannot \texttt{delete} a popped node 
          while another thread might still be reading \texttt{old->next}. Real systems 
          use hazard pointers, epoch-based reclamation, RCU-style schemes, or a 
          freelist with enough tagging --- not raw \texttt{delete} on the hot path.
\end{itemize}
\par\noindent\textbf{Interview tip:} draw the CAS loop on \texttt{head}; name 
\highlight{ABA} and \highlight{reclamation} as the reasons production stacks are 
more than ten lines. HFT often avoids unbounded Treiber stacks on the hot path 
(allocation + reclamation) and prefers preallocated SPSC/MPMC rings.

\section{Lock-Free Queue (Michael--Scott)}
\label{sec:michael-scott-queue}
See~\cite{wiki-michael-scott-queue},~\cite{michael-scott-1996}. The 
\highlight{Michael--Scott} queue is the classic lock-free \highlight{MPMC FIFO}: 
atomic \texttt{head} and \texttt{tail}, a dummy node so enqueue and dequeue touch 
different ends, and CAS on \texttt{next} / \texttt{head} / \texttt{tail}.
\newline
Harder than Treiber; worth knowing the story even if you ship a bounded ring in 
production.

\subsection{Structure}
\label{subsec:ms-structure}
Treiber is a stack (LIFO) with one CAS site. Michael--Scott is a FIFO, so enqueue and 
dequeue must touch \highlight{different} ends without locking --- that is why the 
dummy node exists.
\begin{itemize}
    \item Each node holds a value and an atomic \texttt{next} (initially null).
    \item The queue always contains at least a \highlight{dummy} node. 
          \texttt{head} points at the dummy (dequeue end); \texttt{tail} points at 
          the last node (enqueue end).
    \item \highlightbf{enqueue} swings \texttt{tail->next}, then helps advance 
          \texttt{tail}.
    \item \highlightbf{dequeue} advances \texttt{head} over the dummy to the first 
          real node, returns that value, and makes the old first node the new dummy.
\end{itemize}
\textbf{Noise removed by the dummy:} enqueue and dequeue rarely contend on the same 
pointer in the empty/steady cases; without it, empty-queue races on a single 
\texttt{head}/\texttt{tail} pair become much nastier.

\subsection{Enqueue / Dequeue (logic)}
\label{subsec:ms-enqueue-dequeue}
Helping is the other core idea: if \texttt{tail} lags behind a linked \texttt{next}, 
any thread may CAS \texttt{tail} forward. That keeps the algorithm lock-free when a 
producer stalls mid-enqueue.
\newline
\textbf{Enqueue} of value $v$:
\begin{enumerate}
    \item Allocate node $n$ with value $v$, \texttt{next = nullptr}.
    \item Loop: read \texttt{tail}; read \texttt{tail->next}.
    \item If \texttt{tail} is still the observed tail:
          \begin{itemize}
              \item If \texttt{tail->next} is null, CAS it to $n$ (link). On success, 
                    try to CAS \texttt{tail} to $n$ (may fail --- another thread helps).
              \item Else (tail lagging), CAS \texttt{tail} forward to \texttt{tail->next} 
                    (help) and retry.
          \end{itemize}
\end{enumerate}
\textbf{Dequeue:}
\begin{enumerate}
    \item Loop: read \texttt{head}, \texttt{tail}, and \texttt{head->next} (first real).
    \item If \texttt{head} is still current:
          \begin{itemize}
              \item If \texttt{head == tail} and \texttt{head->next} is null $\Rightarrow$ empty.
              \item If \texttt{head == tail} and \texttt{head->next} is non-null $\Rightarrow$ 
                    help advance \texttt{tail}, retry.
              \item Else CAS \texttt{head} from old head to \texttt{head->next}; on success 
                    the value in the new head (former first real node) is yours; retire 
                    the old dummy.
          \end{itemize}
\end{enumerate}
Helping keeps the algorithm lock-free: a stalled enqueuer does not permanently block 
dequeuers.

\subsection{Implementation Sketch}
\label{subsec:ms-implementation}
Didactic sketch (same ABA / reclamation caveats as Treiber; not production-ready):

\begin{lstlisting}[language=C++]
#include <atomic>
#include <optional>
#include <utility>

template <typename T>
class MichaelScottQueue {
    struct Node {
        std::optional<T> value;                 // dummy starts empty
        std::atomic<Node*> next{nullptr};
        explicit Node(std::optional<T> v = std::nullopt)
            : value(std::move(v)) {}
    };

    std::atomic<Node*> head_;
    std::atomic<Node*> tail_;

public:
    MichaelScottQueue()
    {
        Node* dummy = new Node{};               // no value
        head_.store(dummy, std::memory_order_relaxed);
        tail_.store(dummy, std::memory_order_relaxed);
    }

    void enqueue(T v)
    {
        Node* n = new Node{std::move(v)};
        for (;;) {
            Node* t = tail_.load(std::memory_order_acquire);
            Node* next = t->next.load(std::memory_order_acquire);
            if (t != tail_.load(std::memory_order_acquire)) {
                continue;                       // stale tail
            }
            if (next == nullptr) {
                if (t->next.compare_exchange_weak(
                        next, n,
                        std::memory_order_release,
                        std::memory_order_relaxed)) {
                    tail_.compare_exchange_strong(
                        t, n,
                        std::memory_order_release,
                        std::memory_order_relaxed);  // help OK if fails
                    return;
                }
            } else {
                tail_.compare_exchange_weak(
                    t, next,
                    std::memory_order_release,
                    std::memory_order_relaxed);  // help advance tail
            }
        }
    }

    bool try_dequeue(T& out)
    {
        for (;;) {
            Node* h = head_.load(std::memory_order_acquire);
            Node* t = tail_.load(std::memory_order_acquire);
            Node* next = h->next.load(std::memory_order_acquire);
            if (h != head_.load(std::memory_order_acquire)) {
                continue;
            }
            if (h == t) {
                if (next == nullptr) {
                    return false;               // empty
                }
                tail_.compare_exchange_weak(
                    t, next,
                    std::memory_order_release,
                    std::memory_order_relaxed);
                continue;
            }
            // next is the first real node; claim it by swinging head
            if (head_.compare_exchange_weak(
                    h, next,
                    std::memory_order_release,
                    std::memory_order_relaxed)) {
                out = std::move(*next->value);
                // retire h (old dummy) via a reclamation scheme
                return true;
            }
        }
    }
};
\end{lstlisting}

\subsection{HFT Takeaway}
\label{subsec:ms-hft-takeaway}
\begin{itemize}
    \item Michael--Scott teaches \highlight{CAS + helping} and why MPMC FIFO is subtle.
    \item On latency-critical paths, shops often prefer \highlight{bounded} rings 
          (SPSC, or a known MPMC ring with preallocated slots) over unbounded 
          linked queues: no \texttt{new} on the hot path, simpler failure mode 
          (``full''), friendlier to caches.
    \item Still: if an interviewer asks for ``a lock-free queue,'' outlining 
          Michael--Scott (dummy node, CAS next, help tail, ABA/reclamation) is the 
          expected answer --- then say what you would ship instead and why.
\end{itemize}
\par\noindent\textbf{Interview tip:} SPSC = indices + release/acquire; Treiber = CAS 
\texttt{head}; Michael--Scott = dummy + CAS \texttt{next}/\texttt{head}/\texttt{tail} 
+ helping. Always mention \highlight{ABA} and \highlight{memory reclamation}.

\section{Synchronization}
\label{sec:synchronization-theory}

\subsection{C++ Reference}
\label{subsec:sync-reference}
Again a deliberate split: Chapter~\ref{chap:concurrency}, 
Section~\ref{subsec:mutual-exclusion} lists \texttt{std::mutex}, 
\texttt{std::shared\char`_mutex}, RAII wrappers, and condition variables. 
This chapter asks \highlight{when} those primitives hurt --- contention, convoys, 
and when a spinlock or SPSC handoff is the better latency story.
\par\noindent\textbf{Interview tip:} name the Concurrency type, then say the performance 
failure mode from the subsections below.

\subsection{Spinlock}
\label{subsec:spinlock}
A \highlight{spinlock} waits by spinning on an atomic flag instead of sleeping in the kernel. 
On a \highlight{dedicated core} with very short critical sections, it can beat a mutex 
(no context switch). On a busy machine, spinning wastes CPU and can hurt latency elsewhere.
\newline
Rough rule: mutex for most code; spinlock only when you control the core and the critical 
section is tiny. See also busy polling (Section~\ref{sec:busy-polling}).

\begin{lstlisting}[language=C++]
#include <atomic>

struct Spinlock {
    std::atomic_flag flag = ATOMIC_FLAG_INIT;

    void lock()
    {
        while (flag.test_and_set(std::memory_order_acquire)) {
            // spin
        }
    }

    void unlock()
    {
        flag.clear(std::memory_order_release);
    }
};
\end{lstlisting}

\section{Performance Considerations}
\label{sec:sync-performance}

\subsection{Lock Contention}
\label{subsec:lock-contention}
\highlight{Contention} means many threads fight for the same lock. Hold times grow, cores wait, 
and caches bounce the lock word (coherence traffic --- Chapter~\ref{chap:cacheperformance}).
\newline
\textbf{Noise removed} by shrinking the critical section or removing the share: wait time, 
futex wake storms, and lock-word ping-pong. Fixes: finer-grained locks, shard data, or 
remove sharing (SPSC / thread-local). On a latency path, ``add another mutex'' is rarely 
the first answer.
\par\noindent\textbf{Interview tip:} contention is a \highlight{sharing} problem; name what 
data is shared before naming lock types.

\subsection{Lock Convoy}
\label{subsec:lock-convoy}
A \highlight{lock convoy} happens when threads repeatedly wake, find the lock taken, and sleep 
again --- often with priority or scheduling effects that serialize progress. Symptom: high 
context-switch rate around one mutex without useful work between switches.
\newline
\textbf{Noise:} wake $\rightarrow$ try $\rightarrow$ sleep loops that burn scheduler time 
while the critical section stays long. Prefer shorter critical sections, fewer waiters, 
and less sharing --- not a higher priority on every waiter.

\subsection{Synchronization Scalability}
\label{subsec:scalability}
Ask: does adding cores increase useful work, or only increase waiting on one lock? 
Designs that scale typically \highlight{partition} work (per-core queues, sharding) instead of 
one global mutex.
\newline
\textbf{Noise removed by partitioning:} serialising the whole machine on one lock word. 
HFT systems often accept \highlight{single-threaded} hot paths plus SPSC handoffs rather than 
fine-grained locking everywhere --- scalability by \emph{not} sharing the hot data structure.
\par\noindent\textbf{Interview tip:} ``more threads'' is not scalability if they all queue on 
one mutex; say how work is partitioned.
